\chapter{10\,000 तक की संख्याएँ}\label{ch:g3:numbers}

एक हज़ार, दो हज़ार, तीन हज़ार\,\dots\ बड़ी संख्याएँ छोटी
संख्याओं से बनती हैं, एक शानदार विचार की मदद से: किसी अंक का
\emph{स्थान} ही उसका मान बताता है। यह अध्याय $10\,000$ तक की
संख्याएँ पढ़ना, लिखना, तुलना करना और क्रम में लगाना सिखाता है।

\section{संख्याएँ पढ़ना और लिखना}

\begin{definition}[स्थानीय मान]\label{def:g3:numbers:place}
संख्या $0, 1, 2, \dots, 9$ अंकों से लिखी जाती है। दाईं ओर से
पढ़ने पर स्थान इस क्रम में आते हैं: इकाई, दहाई, सैकड़ा, हज़ार।
$3\,254$ में:
\[
3\,254 = 3 \text{ हज़ार} + 2 \text{ सैकड़ा} + 5 \text{ दहाई}
+ 4 \text{ इकाई}.
\]
\end{definition}

\begin{omfigure}
\begin{tikzpicture}[scale=0.9]
  \draw (0,0) grid[xstep=2.3, ystep=0.85] (9.2,1.7);
  \node[font=\small] at (1.15,1.275) {हज़ार};
  \node[font=\small] at (3.45,1.275) {सैकड़ा};
  \node[font=\small] at (5.75,1.275) {दहाई};
  \node[font=\small] at (8.05,1.275) {इकाई};
  \node[font=\large, omDef] at (1.15,0.425) {$3$};
  \node[font=\large, omDef] at (3.45,0.425) {$2$};
  \node[font=\large, omDef] at (5.75,0.425) {$5$};
  \node[font=\large, omDef] at (8.05,0.425) {$4$};
\end{tikzpicture}

{\small $3\,254$ की स्थानीय-मान सारणी; इसे ``तीन हज़ार दो सौ
चौवन'' पढ़ते हैं।}
\end{omfigure}

\begin{example}[स्थान रोकने वाला शून्य]\label{ex:g3:numbers:zero}
``चार हज़ार सात'' को $4\,007$ लिखा जाता है: हज़ार में $4$,
इकाई में $7$, और बीच में \emph{\omterm{def:g1:counting:zero}{शून्य}}, ताकि हर अंक अपने सही
स्थान पर टिका रहे। उनके बिना $47$ बिलकुल दूसरी संख्या होती!
\end{example}

\begin{example}[अदला-बदली]\label{ex:g3:numbers:exchange}
दस इकाई मिलकर एक दहाई; दस दहाई मिलकर एक सैकड़ा; दस सैकड़ा मिलकर
एक हज़ार। तो $10$ सैकड़ा $= 1\,000$, और $32$ दहाई
$= 320$: संख्या $320$ में $32$ दहाइयाँ हैं (केवल अंक
$2$ नहीं!)।
\end{example}

\section{तुलना और क्रम}

\begin{method}[दो संख्याओं की तुलना]\label{met:g3:numbers:compare}
\begin{enumerate}
  \item अंक गिनो: ज़्यादा अंक यानी बड़ी संख्या
        ($1\,002 > 998$)।
  \item अंक बराबर हों: \emph{बाईं} ओर से अंक-दर-अंक तुलना करो।
        पहला \omterm{ex:g1:subtraction:difference}{अंतर} ही फ़ैसला कर देता है:
        $5\,614 > 5\,589$, क्योंकि सैकड़े पर $6 > 5$।
  \item उत्तर $<$ (से छोटी) या $>$ (से बड़ी) चिह्न से लिखो:
        खुला मुँह हमेशा बड़ी संख्या की ओर रहता है।
\end{enumerate}
\end{method}

\begin{omfigure}
\begin{tikzpicture}[scale=1.1]
  \draw[-Stealth] (-0.3,0) -- (10.5,0);
  \foreach \i in {0,1,...,10} \draw (\i,-0.08) -- (\i,0.08);
  \node[below, font=\small] at (0,-0.1) {$0$};
  \node[below, font=\small] at (5,-0.1) {$500$};
  \node[below, font=\small] at (10,-0.1) {$1000$};
  \fill[omThm] (2.4,0) circle (2pt) node[above, font=\small] {$240$};
  \fill[omDef] (6.9,0) circle (2pt) node[above, font=\small] {$690$};
  \fill[omMeth] (9.05,0) circle (2pt) node[above, font=\small] {$905$};
\end{tikzpicture}

{\small सैकड़ों में अंकित रेखा पर संख्याएँ: जितना दाईं ओर, उतनी
बड़ी। $240 < 690 < 905$।}
\end{omfigure}

\begin{example}\label{ex:g3:numbers:order}
सबसे छोटी से सबसे बड़ी के क्रम में लगाओ: $780$, $8\,700$, $807$, $87$।
पहले अंकों की संख्या से: $87$ (दो अंक), फिर $780$ और $807$
(तीन), फिर $8\,700$ (चार)। $780$ और $807$ में: सैकड़े के अंक
देखो, $7 < 8$, इसलिए $780 < 807$। अंतिम क्रम:
\[
87 < 780 < 807 < 8\,700 .
\]
\end{example}

\section{सम और विषम}

\begin{definition}[सम, विषम]\label{def:g3:numbers:evenodd}
कोई संख्या \emph{सम संख्या}\index{सम संख्या} तब कहलाती है जब उसे
दो बराबर पूरे हिस्सों में बाँटा जा सके --- उसका इकाई अंक $0$, $2$,
$4$, $6$ या $8$ होता है। नहीं तो वह
\emph{विषम संख्या}\index{विषम संख्या} कहलाती है (इकाई अंक $1$,
$3$, $5$, $7$ या $9$)।
\end{definition}

\begin{example}
$46$ सम है ($46 = 23 + 23$); $47$ विषम है (दो बराबर हिस्सों के
लिए आधी वस्तु चाहिए होती)। केवल \emph{आख़िरी} अंक मायने रखता है:
$3\,578$ सम है, $2\,341$ विषम।
\end{example}

\begin{omfigure}
\begin{tikzpicture}[scale=0.45]
  \foreach \i in {0,...,5} \fill[omDef] ({\i},0) circle (0.28);
  \foreach \i in {0,...,5} \fill[omDef] ({\i},1) circle (0.28);
  \node[below, font=\small] at (2.5,-0.7) {$12$: सम, दो बराबर पंक्तियाँ};
  \begin{scope}[xshift=10.5cm]
  \foreach \i in {0,...,4} \fill[omThm] ({\i},0) circle (0.28);
  \foreach \i in {0,...,4} \fill[omThm] ({\i},1) circle (0.28);
  \fill[omThm] (5,0) circle (0.28);
  \draw[omExo, thick] (5,0) circle (0.5);
  \node[below, font=\small] at (2.5,-0.7) {$11$: विषम, एक बिंदु अकेला};
  \end{scope}
\end{tikzpicture}

{\small \omterm{def:g3:numbers:evenodd}{सम संख्याएँ} दो बराबर पंक्तियों में बँट जाती हैं; \omterm{def:g3:numbers:evenodd}{विषम
संख्याएँ} हमेशा एक को अकेला छोड़ देती हैं।}
\end{omfigure}

\section{अभ्यास}

\begin{exercise}[$\star$]\label{exo:g3:numbers:1}
अंकों में लिखो: ``दो हज़ार तीन सौ पैंतालीस''; ``छह
हज़ार पचास''; ``नौ हज़ार एक''।
\end{exercise}

\begin{exercise}[$\star$]\label{exo:g3:numbers:2}
शब्दों में लिखो: $1\,208$; \quad $4\,070$; \quad $9\,999$।
\end{exercise}

\begin{exercise}[$\star$]\label{exo:g3:numbers:3}
$7\,382$ में: दहाई के स्थान पर कौन-सा अंक है? हज़ार के स्थान
पर? संख्या में कुल कितने \emph{सैकड़े} हैं
(सावधान: केवल सैकड़े का अंक नहीं ---
\cref{ex:g3:numbers:exchange} याद करो)?
\end{exercise}

\begin{exercise}[$\star$]\label{exo:g3:numbers:4}
\cref{def:g3:numbers:place} की तरह हिस्सों में बाँटो:
$5\,631$; \quad $2\,046$; \quad $8\,500$।
\end{exercise}

\begin{exercise}[$\star$]\label{exo:g3:numbers:5}
उतार कर $<$ या $>$ से पूरा करो:
\[
456 \;?\; 465, \qquad
1\,203 \;?\; 989, \qquad
6\,090 \;?\; 6\,009, \qquad
9\,999 \;?\; 10\,000 .
\]
\end{exercise}

\begin{exercise}[$\star$]\label{exo:g3:numbers:6}
सबसे छोटी से सबसे बड़ी के क्रम में लिखो: $530$; \ $3\,500$; \ $353$; \ $5\,033$;
\ $503$।
\end{exercise}

\begin{exercise}[$\star$]\label{exo:g3:numbers:7}
$0$ से $1\,000$ तक सैकड़ों में अंकित संख्या रेखा पर $150$, $480$,
$820$ और $990$ को (लगभग) रखो।
\end{exercise}

\begin{exercise}[$\star$]\label{exo:g3:numbers:8}
सम या विषम? \ $34$; \ $57$; \ $70$; \ $2\,481$; \ $6\,548$।
\end{exercise}

\begin{exercise}[$\star$]\label{exo:g3:numbers:9}
कदम-दर-कदम गिनो: $2\,970$ से $10$ के पाँच कदम; $860$ से
$100$ के चार कदम; $9\,996$ से $1$ के चार कदम।
आख़िरी वाले के अंत में क्या दिखता है?
\end{exercise}

\begin{exercise}[$\star\star$]\label{exo:g3:numbers:10}
अंकों $3$, $7$, $0$, $5$ में से हर एक को ठीक एक बार लेकर चार
अंकों की सबसे बड़ी संख्या लिखो, फिर सबसे छोटी (संख्या $0$ से
शुरू नहीं हो सकती)।
\end{exercise}

\begin{exercise}[$\star\star$]\label{exo:g3:numbers:11}
मैं तीन अंकों की संख्या हूँ। मेरे अंक $2$, $5$ और $8$ हैं, मैं सम हूँ,
और $800$ से बड़ी हूँ। मैं कौन हूँ?
\end{exercise}
